\documentclass[12pt]{article}
\usepackage[T1]{fontenc}
\usepackage[utf8]{inputenc}
\usepackage{lmodern}
\usepackage{textcomp}
\usepackage{enumitem}
\usepackage{geometry}
\usepackage{graphicx}
\usepackage{amsmath, amssymb}
\geometry{margin=1in}
\begin{document}
\begin{center}
History - K-12 Level Assessment 14 \
Total Marks: 40 \
Answer all questions.
\end{center}

\section*{Multiple Choice (10 marks)}
\begin{enumerate}[label=\arabic*.]
\item Which commonly used name for a culture was borrowed from another Native American language and literally means "ancient enemies"?
\begin{enumerate}[label=(\alph*)]
    \item Kwakiutl
    \item Hohokam
    \item Anasazi
    \item Mogollon
\end{enumerate}
\item The large, extinct herbivores of the Pleistocene era are referred to as:
\begin{enumerate}[label=(\alph*)]
    \item megaliths.
    \item megafauna.
    \item megaflora.
    \item megatrons.
\end{enumerate}
\item Some of the oldest painted art in the world was recently discovered in caves on the island of Sulawesi in Indonesia. What do the images include and to when do they date?
\begin{enumerate}[label=(\alph*)]
    \item animals and human bodies from more than 15,000 to about 20,000 years ago
    \item animals and hand stencils from more than 35,000 to about 40,000 years ago
    \item animals and hand stencils from more than 45,000 to about 100,000 years ago
    \item animals and abstract geometric shapes from 15,000 to more than 100,000 years ago
\end{enumerate}
\item There is evidence that premodern Homo sapiens existed:
\begin{enumerate}[label=(\alph*)]
    \item 800,000 to 200,000 years ago.
    \item 600,000 to 30,000 years ago.
    \item 400,000 to 10,000 years ago.
    \item 100,000 to 10,000 years ago.
\end{enumerate}
\item During the Pleistocene era, the islands of Java, Sumatra, Bali, and Borneo formed \_\_\_\_\_\_\_\_\_\_\_. The landmass connecting Australia, New Guinea, and Tasmania was called \_\_\_\_\_\_\_\_\_\_\_.
\begin{enumerate}[label=(\alph*)]
    \item Sunda; Wallacea
    \item Wallacea; Beringia
    \item Beringia; Sahul
    \item Sunda; Sahul
\end{enumerate}
\end{enumerate}

\section*{True / False (10 marks)}
\textit{Answer True or False}
\begin{enumerate}[label=\arabic*.]
\setcounter{enumi}{5}
\item True or False: The first metals in South America were mainly produced through casting methods.
\item True or False: Inca architecture is renowned for its precise jigsaw-like stone masonry that fits stones together without the use of mortar.
\item True or False: Levallois technology marked a shift from producing core tools to producing flakes for more refined toolmaking.
\item True or False: Genes inherited from Neandertals are associated with the metabolism of fat and certain human diseases like Lupus in modern Europeans.
\item True or False: The Shang Dynasty laid the foundation for a coherent development of Chinese civilization that lasted into the 13 th century.
\end{enumerate}

\section*{Long Form Response (20 marks)}
\textit{Answer each question with a clear and complete explanation.}
\begin{enumerate}[label=\arabic*.]
\setcounter{enumi}{10}
\item Discuss the various methods historically used to remove soil from around site materials, and analyze their significance in archaeological practices and site preservation.
\item Discuss the differences between absolute and relative dating in historical context, providing examples of each and explaining their significance in understanding timelines.
\end{enumerate}

\vfill
\noindent\textit{End of Paper}
\end{document}